\section{Dimension reduction for the limiting measures}
\label{sec:dimension-reduction}

We now separate the final argument from the varifold construction.
The preceding sections produced a nonzero height measure, its flux
and stress measures, and a density inequality at many centers. The
following proposition shows that these data cannot have more than
$m-1$ dimensions of such centers. Its proof uses a second blow-up:
the frequency becomes constant at several independent centers, and
equality in a positive matrix measure then forces the height to vanish.

Fix an integer $p\ge4$. For $x\in\R^m$ and $l>0$ write
\begin{equation}
 \label{eq:reduction-kernels}
 \phi_{x,l}(y)=l^{-m}\left(1-\frac{|y-x|^2}{l^2}\right)_+^p,
 \qquad
 \psi_{x,l}(y)=2p\,l^{-m}
       \left(1-\frac{|y-x|^2}{l^2}\right)_+^{p-1}.
\end{equation}
We retain the kernels of \eqref{eq:limit-kernels} in order to state
the measure-theoretic argument independently.

\begin{proposition}[Dimension reduction for height and stress]
\label{prop:measure-reduction}
Let $\Omega\subset\R^m$ be connected and open. Let $\nu$ be a
nonzero nonnegative Radon measure, $f$ an $\R^m$-valued Radon measure,
$\beta$ a symmetric matrix-valued Radon measure, and $\mathfrak m$
a signed Radon measure on $\Omega$. Suppose that
\begin{equation}
 \label{eq:abstract-psd}
 \begin{pmatrix}\nu&f^{\rm t}\\ f&\beta\end{pmatrix}
 \quad\hbox{is a positive semidefinite matrix-valued measure},
 \qquad d_0=\tr\beta,
\end{equation}
and that, in distributions,
\begin{equation}
 \label{eq:abstract-identities}
 \nabla\nu=2f,\qquad \diver f=d_0,\qquad
 \int\diver X\,\dd\mathfrak m
       =\int DX:\dd\beta
 \quad(X\in C_c^\infty(\Omega;\R^m)).
\end{equation}
Assume also that for every open $\Omega'\Subset\Omega$ there is a
constant $C_{\Omega'}$ such that, on $\Omega'$,
\begin{equation}
 \label{eq:abstract-sign}
 2\mathfrak m\le d_0,\qquad |\mathfrak m|\le C_{\Omega'}d_0.
\end{equation}
Let $Z\subset\Omega$ be a Borel set such that for each $x\in Z$,
\begin{equation}
 \label{eq:abstract-center}
 2l^2\int\phi_{x,l}\,\dd\mathfrak m
       \ge\int\psi_{x,l}\,\dd\nu
\end{equation}
for every sufficiently small $l>0$. Then
$\dim_{\Hh}Z\le m-1$.
\end{proposition}

All tensor inequalities in this statement are measure inequalities.
For instance, \eqref{eq:abstract-psd} means that contraction with
every constant vector in $\R^{m+1}$ gives a nonnegative measure.
Equivalently, the matrix density with respect to its trace measure is
positive semidefinite almost everywhere. This latter formulation also
allows contraction with bounded, point-dependent vectors.

\subsection{Kernel identities and positive height}

For the moment fix an arbitrary center $x\in\Omega$, not necessarily
in $Z$, and restrict to radii with
$\overline B_l(x)\subset\Omega$. Define
\begin{equation}
 \label{eq:abstract-frequencies}
 \begin{aligned}
 H(l)&=\int\psi_{x,l}\,\dd\nu,&
 A(l)&=l^2\int\phi_{x,l}\,\dd\mathfrak m,\\
 D(l)&=l^2\int\phi_{x,l}\,\dd d_0,&
 R(l)&=\int\psi_{x,l}(y)\,\dd\beta(y)[y-x,y-x].
 \end{aligned}
\end{equation}
Dots denote derivatives with respect to $\log l$. The distributional
identities imply
\begin{equation}
 \label{eq:exact-frequency-identities}
 \dot A=2A-D+R,\qquad
 D=\int\psi_{x,l}(y)(y-x)\cdot\dd f(y),\qquad
 \dot H=2D,\qquad RH\ge D^2,\qquad 2A\le D.
\end{equation}
Here is the calculation. The identities
\[
 \nabla\phi_{x,l}=-l^{-2}(y-x)\psi_{x,l},\qquad
 \dot\psi_{x,l}=-\diver\big((y-x)\psi_{x,l}\big)
\]
give the second and third formulas. In the stress identity use
$X=(y-x)\phi_{x,l}$; differentiating $A$ then gives the first
formula. Cauchy--Schwarz for the positive matrix measure gives
$D^2\le RH$, and the signed inequality gives $2A\le D$.
Although the kernels are compactly supported powers rather than
$C^\infty$ functions, their use is justified by approximation in
$C^2$ on a compact interior set. The choice $p\ge4$ provides that
regularity. All quantities and their displayed derivatives are
continuous in the admissible center and radius.

We need positive height before dividing by $H$. The next lemma is
the finite-radius continuation argument used in the failed-fit
analysis of AE; we give it in the measure form required here
\cite[Proof of Lemma~7.2]{AE}.

\begin{lemma}[No open zero-height region]
\label{lem:positive-limit-height}
For the nonzero measure $\nu$ on the connected set $\Omega$ in
Proposition~\ref{prop:measure-reduction}, equations
\eqref{eq:abstract-psd}--\eqref{eq:abstract-sign} imply $H(l)>0$
whenever $\overline B_l(x)\subset\Omega$ and $l>0$.
\end{lemma}

\begin{proof}
Suppose that $\nu$ vanishes on a nonempty open subset. Because
$\nu\ne0$ and $\Omega$ is connected, choose a point $x$ outside
$\spt\nu$ so close to its boundary that
\[
 l_0=\operatorname{dist}(x,\spt\nu)>0,
 \qquad \overline B_{2l_0}(x)\subset\Omega.
\]
Then $H(l)=0$ for $l\le l_0$ and $H(l)>0$ for
$l_0<l\le l_1$, where $l_1>l_0$ is fixed and interior. Put
$t=\log l$, $t_j=\log l_j$, and on $(t_0,t_1]$ define
\[
 P=A/H,\qquad n=D/H,\qquad q=R/H.
\]
By \eqref{eq:exact-frequency-identities},
\begin{equation}
 \label{eq:finite-radius-frequency}
 n\ge0,\quad n\ge2P,\quad q\ge n^2,\qquad
 \dot P=2P-n+q-2Pn
       \ge(n-1)(n-2P).
\end{equation}
For $P\ge1$ the right-hand side is nonnegative. Therefore, going
inward from $t_1$,
$P(t)\le\max\{1,P(t_1)\}$.

On $B_{l_0}(x)$, the distributional identities give $f=0$ and
$d_0=0$, and \eqref{eq:abstract-sign} gives $\mathfrak m=0$.
The compact kernels converge uniformly and vanish on the boundary
sphere, so $A(t)\to0$ as $t\downarrow t_0$, even if
$\mathfrak m$ charges that sphere. Moreover,
\[
 (e^{-2t}A)^{\displaystyle\cdot}
   =e^{-2t}(R-D)\ge-\tfrac14e^{-2t}H,
\]
because $q-n\ge n^2-n\ge-1/4$. Since $\dot H=2D\ge0$,
integration from $t_0$ yields
\[
 P(t)\ge-\frac{e^{2(t-t_0)}-1}{8}.
\]
Thus $P$ is bounded on this finite interval. Using $q\ge n^2$ in
the exact equation for $\dot P$ now gives
\[
 \dot P\ge\tfrac12 n^2-C.
\]
Integrating from $t_0+\varepsilon$ to $t_1$ bounds
$\int_{t_0+\varepsilon}^{t_1}n^2\,\dd t$ independently of
$\varepsilon$. In particular $n$ is integrable down to $t_0$.
But $(\log H)^{\displaystyle\cdot}=2n$, which prevents $H$ from
tending to zero there. This contradiction excludes every nonempty
open zero-height region. Since the kernel is positive in its open
ball, the stated conclusion follows.
\end{proof}

\subsection{A second blow-up at the centers}

\begin{proof}[Proof of Proposition~\ref{prop:measure-reduction}]
The preceding lemma makes $P(x,l)=A(x,l)/H(x,l)$ well defined at
every interior ball. In the harmonic-map model from
Section~\ref{sec:linearized}, this quotient is half the weighted
Dirichlet frequency
\[
 \frac{l^2\int\phi_{x,l}|Du|^2\,\dd y}
      {\int\psi_{x,l}|u|^2\,\dd y}.
\]
For a nonzero harmonic map homogeneous of degree $d$ about $x$, the
flux identity in \eqref{eq:exact-frequency-identities} makes this
frequency equal to $d$, so $P=d/2$. For the measure data, constant
frequency at the retained centers, together with the lower bound
$2P\ge1$ supplied by their density comparison, will force the radial
measure identity used below.

At a center in $Z$, the additional hypothesis
\eqref{eq:abstract-center} gives $2P\ge1$ for small radii. Hence
\eqref{eq:finite-radius-frequency} improves to
\begin{equation}
 \label{eq:center-frequency-monotonicity}
 2P\ge1,\qquad n\ge2P,\qquad
 \dot P\ge(n-1)(n-2P)\ge0.
\end{equation}
In particular the finite limit $P(x,0+)$ exists and is at least
$1/2$.

Suppose that $\dim_{\Hh}Z>m-1$, and choose
$m-1<s'<\dim_{\Hh}Z$. We may first restrict to centers in a compact
interior region for which \eqref{eq:abstract-center} holds on one
common interval $0<l<l_*$. To justify this restriction, exhaust
$\Omega$ by compact sets with positive boundary distance and impose
the comparison for all rational radii below $1/k$. The resulting
sets are Borel by continuity of the kernels, their union is $Z$, and
continuity in radius supplies the comparison at every radius. One
of them still has dimension greater than $s'$.

The compact-subset theorem and Egoroff's theorem now give a compact
set $K$ of positive finite $\Hh^{s'}$ measure on which
\begin{equation}
 \label{eq:uniform-center-frequency}
 P(y,l)\longrightarrow P(y,0+)\quad\hbox{uniformly as }l\downarrow0.
\end{equation}
Indeed first apply Egoroff at dyadic radii; monotonicity supplies the
same conclusion between them. Since each function $P(\cdot,l)$ is
continuous, its uniform limit is continuous on $K$. A further
compact restriction gives
\begin{equation}
 \label{eq:second-center-frostman}
 \Hh^{s'}(K\cap B_u(y))\le C u^{s'}.
\end{equation}
Choose $x\in K$ with positive upper $s'$-density, and radii
$l_j\downarrow0$ such that
$\Hh^{s'}(K\cap B_{l_j}(x))\ge c l_j^{s'}$.

We next normalize the measure tuple at these scales. Fix an open
neighborhood $\Omega'\Subset\Omega$ of $x$, so that the constant
in \eqref{eq:abstract-sign} is fixed on $\Omega'$. Put
$T_j(y)=(y-x)/l_j$ and $N_j=l_j^m H(x,l_j)$. On the expanding
domains $T_j(\Omega)$ define
\begin{equation}
 \label{eq:second-measure-normalization}
 \nu_j=N_j^{-1}(T_j)_\#\nu,\qquad
 f_j=l_jN_j^{-1}(T_j)_\# f,\qquad
 (\beta_j,d_j,\mathfrak m_j)
   =l_j^2N_j^{-1}(T_j)_\#(\beta,d_0,\mathfrak m).
\end{equation}
The identities and inequalities
\eqref{eq:abstract-psd}--\eqref{eq:abstract-sign} are invariant under
this normalization, and the new height at center zero and radius one
equals one.

Here are the local mass bounds needed to pass to a limit. At the
fixed center $x$, $P$ is bounded on $0<l<l_*/2$. The exact quotient
equation and $n\le\sqrt q$ therefore give
\[
 \dot P=q-(1+2P)n+2P\ge\tfrac12q-C.
\]
On every interval of radii with bounded ratio, this bounds
$\int q\,\dd\log l$. Since
$\dot H/H=2n\le2\sqrt q$, Cauchy--Schwarz gives two-sided
comparison of $H$ on the same interval.

To see the resulting mass estimates explicitly, fix $R>0$ and discard
finitely many $j$ so that all balls used below lie in $\Omega'$.
Comparison on the interval from $\min\{1,R\}l_j$ to
$\max\{1,4R\}l_j$ gives
$H(x,\tau l_j)\le C_RH(x,l_j)$ for $R\le\tau\le4R$.
The kernel $\psi_{x,2Rl_j}$ is at least $c_Rl_j^{-m}$ on
$B_{Rl_j}(x)$; hence
\[
 \nu(B_{Rl_j}(x))\le C_Rl_j^mH(x,l_j).
\]
The integrated bound for $q$ on $[2Rl_j,4Rl_j]$ supplies
$\tau_j\in[2R,4R]$ with $n(x,\tau_jl_j)\le C_R$.
The kernel $\phi_{x,\tau_jl_j}$ is also at least
$c_Rl_j^{-m}$ on $B_{Rl_j}(x)$, so the definition of $D=nH$ gives
\[
 d_0(B_{Rl_j}(x))\le C_Rl_j^{m-2}H(x,l_j),\qquad
 \nu_j(B_R)+d_j(B_R)\le C_R.
\]
Positivity bounds $\beta_j$ by its trace $d_j$, and
Cauchy--Schwarz for \eqref{eq:abstract-psd} bounds $f_j$.
Finally, $|\mathfrak m_j|\le C_{\Omega'}d_j$ on every fixed ball
for all sufficiently large $j$. In particular the same constant
is available on all fixed balls of the eventual limit.

Pass to a diagonal subsequence with local weak limits on $\R^m$,
and denote them by
\[
 \nu_\infty,\quad f_\infty,\quad\beta_\infty,\quad
 d_\infty,\quad\mathfrak m_\infty.
\]
They satisfy the same measure identities and
inequalities. The normalization gives
$H_\infty(0,1)=1$, so $\nu_\infty\ne0$ and
Lemma~\ref{lem:positive-limit-height} applies on all of $\R^m$.

The rescaled center measures
\[
 \lambda_j=l_j^{-s'}(T_j)_\#
       \big(\Hh^{s'}\!\llcorner(K\cap\overline B_{l_j}(x))\big)
\]
have a nonzero weak limit $\lambda$ on $\overline B_1$, with the
upper ball bound \eqref{eq:second-center-frostman}. Consequently
$\spt\lambda$ spans $\R^m$. Otherwise its bounded support lies in
an $(m-1)$-dimensional linear space and can be covered by
$O(\varepsilon^{-(m-1)})$ balls of radius $\varepsilon$. The ball
bound would give
$\lambda(\R^m)\le C\varepsilon^{s'-m+1}\to0$, contradicting
its nonzero mass.

We have now retained enough centers in the second blow-up. Its
remaining feature is that they all have the same constant frequency.
If $z\in\spt\lambda$, choose $y_j\in K$ with
$T_j(y_j)\to z$. At every fixed radius $r>0$,
\[
 P_j(T_j(y_j),r)=P(y_j,l_jr)
       \longrightarrow P(x,0+)=:P_0.
\]
This follows from \eqref{eq:uniform-center-frequency} and continuity
of the limiting frequency on $K$. Local mass bounds and continuous
compact kernels permit passage to the weak limits even with these
moving centers. Since the limiting height is positive,
\begin{equation}
 \label{eq:constant-frequency-centers}
 P_\infty(z,r)=P_0
 \quad(z\in\spt\lambda\cup\{0\},\ r>0),
 \qquad P_0\ge\tfrac12.
\end{equation}
The comparison inequality also passes to these centers.

\subsubsection*{Equality at independent centers}

Fix $z\in\spt\lambda\cup\{0\}$. In the limit,
\eqref{eq:center-frequency-monotonicity} and
\eqref{eq:constant-frequency-centers} give
\[
 0=\dot P_\infty\ge(n-1)(n-2P_0)\ge0.
\]
As $n\ge2P_0\ge1$, this forces $n=2P_0=:n_0$, including the
case $P_0=1/2$. The exact quotient equation then gives
$q=n_0^2$, or equivalently $D=n_0H$ and $R=n_0^2H$.

Let $\mathcal M_\infty$ be the positive matrix measure in
\eqref{eq:abstract-psd} for the limiting data. Contract it with
$(-n_0,y-z)$ and integrate against $\psi_{z,r}$. The result is
\[
 n_0^2H-2n_0D+R=0.
\]
With respect to the trace measure, a positive semidefinite matrix
whose quadratic form vanishes on a vector annihilates that vector.
Since $\psi_{z,r}>0$ on $B_r(z)$, the first row therefore gives
\[
 (y-z)\cdot\dd f_\infty(y)=n_0\,\dd\nu_\infty(y)
 \quad\hbox{on }B_r(z).
\]
The radius is arbitrary, so this is an identity on $\R^m$.
Subtracting the identity for center zero yields
$z\cdot f_\infty=0$. The support of $\lambda$ spans $\R^m$,
and hence $f_\infty=0$. The identity at zero then gives
$n_0\nu_\infty=0$, contradicting $n_0\ge1$ and
$H_\infty(0,1)=1$. This proves the proposition.
\end{proof}

\subsection{Completion of the varifold argument}

\begin{proof}[Proof of Theorem~\ref{thm:main}]
Suppose that $\dim_{\Hh}\Sing(V)>m-1$, and choose
$m-1<s<\dim_{\Hh}\Sing(V)$. Proposition~\ref{prop:flat-scales}
gives the fixed singular point, flat target scales, and center
measures satisfying \eqref{eq:center-measure-bounds}.
Proposition~\ref{prop:centered-blowup} supplies the centered,
nonzero height normalization at these scales.
Proposition~\ref{prop:linearized} then gives on the connected base
ball $B_4\subset\Pi$ the nonzero measure tuple satisfying
\eqref{eq:limit-psd}, \eqref{eq:limit-identities}, and
\eqref{eq:limit-sign}. These are precisely the hypotheses
\eqref{eq:abstract-psd}--\eqref{eq:abstract-sign}.

Pass to a weak limit $\sigma$ of the center measures. It is nonzero,
is supported in $\Pi\cap\overline{\mathbf B}_1$, and retains the
upper ball bound with exponent $s$. The mass distribution principle
therefore gives
$\dim_{\Hh}\spt\sigma\ge s$.
Proposition~\ref{prop:density-comparison} supplies
\eqref{eq:center-comparison} at every point of $\spt\sigma$,
which is \eqref{eq:abstract-center}. Applying
Proposition~\ref{prop:measure-reduction} with
$Z=\spt\sigma$ gives
\[
 s\le\dim_{\Hh}\spt\sigma\le m-1,
\]
a contradiction.

For sharpness, take two distinct $m$-planes intersecting in an
$(m-1)$-plane, each with multiplicity one. Their sum is stationary
by additivity of first variation. Away from the intersection it is
locally a single smooth plane, whereas at the intersection its
support is not a smooth embedded $m$-submanifold. Its singular set
therefore has dimension exactly $m-1$.
\end{proof}
